\begin{explanationbox}
	\textbf{\textcolor{CExplanation}{为什么会想到这些公式？}}

	求 $\log_2(8\times4)$，可以先算 $8\times4=32$，再找 $2$ 的几次方等于 $32$。有没有更快的办法？\textbf{先看 $8$ 和 $4$ 分别是 $2$ 的几次方。}
	\[
		8=2^3,\quad 4=2^2,\quad 8\times4=2^{3+2}=2^5.
	\]
	于是 $\log_2(8\times4)=3+2=5$。\textbf{对数就是在问：“底数的几次方，等于这个数？”}

	\medskip
	\textbf{\textcolor{CExplanation}{把四条规律与旧知识连起来}}
	\begin{description}[style=nextline, leftmargin=2.8em, labelsep=0.5em, itemsep=0.45em]
		\item[\textbf{乘法为什么变加法？}] 同底幂相乘，指数相加：$2^3\times2^2=2^{3+2}$。所以 $\log_2(8\times4)=\log_2 8+\log_2 4$。
		\item[\textbf{除法为什么变减法？}] 同底幂相除，指数相减：$2^3\div2^2=2^{3-2}$。所以 $\log_2(8\div4)=3-2$。
		\item[\textbf{指数为什么能提前？}] 幂再乘方，指数相乘：$(2^3)^2=2^{3\times2}$。所以 $\log_2(8^2)=2\log_2 8$。
		\item[\textbf{为什么能换底？}] 求 $\log_4 8$，就是求 $4^x=8$ 中的 $x$。因为 $4=2^2$、$8=2^3$，所以 $2^{2x}=2^3$，得到 $x=\frac32$。换底公式就是把两边都改写成同一个底数的做法。
	\end{description}

	\medskip
	\textbf{\textcolor{CExplanation}{考试中什么时候用最快？}}
	\begin{itemize}[leftmargin=*, itemsep=0.45em, topsep=0.35em]
		\item \textbf{同底对数有加减：} 优先合并成一个对数，再计算，如 $\log_3 9+\log_3 3=\log_3 27=3$。
		\item \textbf{真数有乘方或根号：} 先把指数提出，如 $\log_2\sqrt8=\frac12\log_2 8$。
		\item \textbf{不同底数相乘：} 先换底，找能约去的部分，如 $\log_2 3\cdot\log_3 4=\log_2 4=2$。
	\end{itemize}

	\textbf{\textcolor{CExplanation}{一句话记住}}

	\textbf{对数没有发明一套新规律，它只是把熟悉的指数运算法则换了一种写法。}
\end{explanationbox}
